\section*{Interpretation of the supplementary analyses}

The broader research question also asks whether sensitive information is used with permission and whether consequential actions receive effective oversight. The graph representation makes specified permissions and review paths testable. The statistical construction evaluates joint risk and workload on a frozen grid. The secondary full-data experiment produces a negative answer for the tested first-alert family: the declared incident-risk target is not achieved, despite useful discrimination and positive challenge utility. Actual clinician effectiveness and patient benefit were not measured. Risk tolerances and review budgets are technical targets, not ethical thresholds derived from patients or clinicians. Changes in simulated consent affect benchmark alert performance and modeled workload; they are not measured causal effects on mortality, treatment or wellbeing.

This is more informative than observing that a bound failed to certify a weak model. The local confidence calculations exclude every tested candidate, and the test frontiers remain far above the 5\% risk target. The calibration pools are large enough that a zero-miss candidate could pass the finite-sample correction. The principal obstacle in this experiment is the combination of score trajectories and the definition of timely first review. More calibration data alone cannot turn a high population-risk candidate into a low-risk one.

The distinction also explains why the stronger predictors do not dominate the linear model on every endpoint. AUROC rewards correct ordering of individual hourly scores. Timely first review depends on the joint trajectory: scores before the window must remain below threshold, while at least one in-window score must cross it. The joint law of the pre-window and in-window maxima is therefore central. An improvement in hourly ranking need not improve that joint event, and a model that sends an early alert can receive utility credit while failing the strict first-review criterion.

Threshold nesting is often invoked to justify risk calibration, but the decision procedure sits between the score and the loss. Alert suppression, cooldown, a one-review limit or a persistent clinical assessment can change the loss family even when the score thresholds are unchanged. The counterexample in the first-alert non-monotonicity proposition isolates this issue without relying on a particular dataset or model architecture.

A useful operational policy may allow an early assessment to remain active, request a second review after deterioration, or suppress repeated notifications for a fixed period. None of these policies is ruled out by our results. Each would need a precise definition of review cost, the duration of an assessment and the event that counts as a successful action. The finite-grid method can evaluate non-monotone policies once their patient-level losses and workloads are specified, although a workload variable above one would require a bounded-mean or count-aware interval rather than the Bernoulli interval used here.

The distinction between infeasible and undetermined should remain visible in such evaluations. A confidence method that only returns accepted candidates cannot, by itself, distinguish an inadequate predictor from an underpowered calibration set. Our lower-bound exclusion criterion supplies that additional information, while limiting the conclusion to the actual finite grid. Extending an exclusion claim to all real-valued thresholds would require additional structure or a uniform result beyond the present analysis.

The graph analysis and queue simulation answer complementary questions. The former asks whether a consequential action can be reached without traversing a review node. The latter asks whether a request receives a timely correct review under explicit assumptions. A path constraint can be satisfied even when every reviewer is overloaded or mistaken; a favorable queue result can coexist with an unmodeled bypass in a real software system.

The queue results are dominated by the upstream incident policy because a reviewer cannot repair a missing or mistimed request within the specified endpoint. This does not imply that reviewer quality is generally unimportant. It shows that improving one component cannot compensate for an event that the system never routes to that component in time. Clinical implementation evidence such as TREWS adoption \citep{trews} is needed to determine how review actually occurs, and externally validated alert-burden analyses such as \citet{wong2021} are needed to assess the workload that a hospital would experience.
